\chapter{Related Work}
\label{ch:RelatedWork}

In model-driven systems, consistency specifications are regarded as the stable component.
It is acknowledged that models may vary across products, the metamodels evolve over time, and existing artifacts are migrated whenever either models or metamodels change.
It is a commonly held assumption that the rules relating to the models remain applicable valid throughout course of the process.
In scenarios where variability is permitted within the rules themselves it is resolved prior to the existence of any artifacts.
Ensuring nothing derived from an earlier variant requires reconsideration.

\section{Variability in model transformations and consistency rules}
\label{sec:RelatedWork:VariabilityInModelTransformationsAndConsistencyRules}

In the context of model transformations research, variability has been identified as a contributing factor, mainly manifesting within the models themselves.
The rules are re-targeted or annotated to process a whole product line, yet they remain fixed.
In instances, where variability is introduced within the rules, it is resolved during the matching process or prior to compilation.
The subsequent papers do not ask of reconsidering the derived transformation.

Strüber et al. have developed a methodology for the integration of similar rules into a unified \texttt{variability-based rule} \cite{struber2018variability}.
Such a rule maintains the union of elements of all the rules it replaces, and parts where rules differ are guarded by presence conditions over variation points.
In the case of the common part, the match is computed only once and specialized per configuration afterward.
The variability is resolved during matching with one artifact representing all variants \cite{struber2018variability}.
The preprocessor described in this thesis resolves annotations earlier, resulting in conventional reactions for the V-SUM without any variability.

The closest architectural precedent extends the \textit{Atlas Transformation Language} (ATL) itself \cite{sijtema2010introducing}.
Sijtema incorporates a \texttt{variability rule} as a construct of the language, which is derived from the standard rule metaclass.
This approach preserves the declarative nature of the original and the implicit scheduling of a standard ATL rule.
Execution semantics are applied precisely when the corresponding variant is selected in the feature model.
It is important to note that this extension does not necessitate a modified engine.
A further transformation compiles the ATL with the new modification into a plain ATL model.
The engine is able to execute those functions, thereby eliminating the need for a plugin \cite{sijtema2010introducing}.
The same reasoning can be applied to annotated reactions.
These reactions are reduced to ordinary reactions prior to compilation, so neither the compiler of the Reactions Language nor Vitruvius needs any support for features.
ATL transformations are unidirectional and stateless.
Consequently, modifications to the variant selection necessitate the implementation of new transformations, which are applied to the source models again \cite{sijtema2010introducing}.
A V-SUM does not exhibit such behavior, given its persistence across iterations.
The paper does not address the artifacts that come from a variant change.
However, it is precisely these artifacts that motivate the migration developed in \autoref{sec:Concept:Migration}.

The expansion of Vitruvius to accommodate variability has been previously explored.
Ananieva et al. propose a unified conceptual model that incorporates variability in space and time.
The interrelated models of a variable system are maintained in a state of consistency as the system evolves.
CPRs form a fixed frame for that construction.
The fundamental system is the variable element in this context.
Consequently, the product line comprises models with a single rule set applicable to all models within the line \cite{ananieva2022preserving}.
In the present work the product line is assembled out of the rules, thereby establishing a complementary relationship between the two.

Neumann et al. integrate the consistency preservation mechanisms of Vitruvius with delta-oriented variability modelling \cite{neumann2025towards}.
A delta consists of operations that add, remove, or modify model elements. Applying a delta transforms one model variant into another.
Deltas are developed for a single domain and applied to a core model that remains consistent across all domains.
Within Vitruvius, applying a delta initiates changes that activate the consistency preservation rules.
The rules identify corresponding changes in other domains, which are then recorded as deltas specific to those domains.
This approach aligns with Vitruvius, as a delta represents a change and Vitruvius maintains consistency by propagating such changes.
Similar to the approach of Ananieva et al., the models exhibit variability, whereas the consistency preservation rules remain unchanged for each variant.
Future work aims to define consistency preservation rules between deltas rather than between metamodels, ensuring that modifications to variability also remain consistent \cite{neumann2025towards}.
This thesis addresses variability in the mapping between metamodels, such as determining whether a UML class is mapped to a Java class or a Java interface.
Such a mapping is defined by the consistency preservation rules, so a delta applied to a model cannot vary it.

\section{Annotative variability}
\label{sec:RelatedWork:AnnotativeVariability}

Annotative variability marks the elements of an artifact directly with the features they belong to \cite{kaestner2008granularity}.
Sulír et al. attach a concern adjacent to the element to which it belongs.
The intent behind a piece of code is embedded within the code itself, rather than being conveyed through a separate document.
Those concerns merely document, rather than derive, as no variant is selected from them.
The mechanism outlined in this thesis employs locality, yet the preprocessor functions by reading reactions and determining which elements get kept \cite{sulir2016recording}.

\section{Migration of existing artifacts}
\label{sec:RelatedWork:MigrationOfExistingArtifacts}

The migration of artifacts whose definition has undergone modifications is a recognized issue, provided that the definition in question is a metamodel.
In the implementation in this thesis, the metamodels are untouched, and the models maintain their validity.
The underlying specifications undergo alterations, resulting in the disruption of correspondence.

In Hearnden et al., the transformation is maintained in a running state and the derivation of each component of the target is recorded \cite{hearnden2006incremental}.
A change can be mapped onto the target incrementally instead of doing everything at once again.
The authors argue that the same mechanism can be applied to a change of the transformation definition itself, rendering it the closest published work to the problem covered in this thesis.
The record is the memory context, while a V-SUM is persisted between runs.
The changed rules have to be recovered afterward from identifiers, hashes and triggers.

Leonhardt et al. commence their investigation by assuming the same observation, relying on changes that have been recorded \cite{leonhardt2015integration}.
A strategy that is employed involves the synthesis of the change sequence that would have generated the models and propagated it afterward.
That represents the internal replay used in \autoref{sec:Concept:Migration}.
Given that the artifacts were built outside Vitruvius, no existing V-SUM, including correspondences is available.
They have to start from an empty V-SUM with a predefined rule \cite{leonhardt2015integration}.
In the thesis correspondences already exist and only the underlying rules change.

\section{Preservation of handwritten content}
\label{sec:RelatedWork:PreservationOfHandwrittenContent}

Greifenberg et al. have conducted a comparative analysis of eight mechanisms that integrate handwritten and generated code, including partial classes and protected regions \cite{greifenberg2015integration}.
Each of them fixes the place where handwritten code may appear in advance.
Vitruvius does not fix such place.
An element no correspondence points at counts as underived.
Everything the rules did not derive ends up in the preservation report (\autoref{sec:Concept:Migration:Preservation}), whether a person wrote it or not.

\section{Other consistency approaches}
\label{sec:RelatedWork:OtherConsistencyApproaches}

Alternative approaches to maintaining consistency in models utilize mechanisms that are not connected to the Reactions language.
Each of these systems operates under the assumption that its own specification is given.
Multi-version models enable the storage of a model's history within a single graph.
This property of multi-version models facilitates transformations that can operate on multiple versions concurrently.
Barkowsky and Giese employ triple graph grammars to define such transformations on multi-version models.
The simultaneous transformation of all versions reduces memory usage and can reduce execution time for extensive version histories.
However, it can add overhead in less favorable circumstances \cite{barkowsky2023incremental}.
It is important to note that reactions utilize the current version of a model exclusively and possess no history.
Graph stores are designed to host versions of the models rather than the rules that relate to them.

König et al. express the correspondence as a query.
The relational operators define views over the models and each of them is consolidated into a triple graph rule.
Those views are analogous to queries in relational databases.
The transformation and synchronization of these views is then managed by the grammars.
In this approach, operators are not regarded as entities that undergo change \cite{konig2025transforming}.

\section{Research gap}
\label{sec:RelatedWork:ResearchGap}

Variability in model transformations exists, but it gets resolved before any artifact is derived from it \cite{struber2018variability,sijtema2010introducing}.
Inside Vitruvius the variability is carried by the models, while the reactions relating them remain fixed \cite{ananieva2022preserving,neumann2025towards}.
The migrating of artifacts whose definition has undergone modification is also well-defined, provided that the definition in question is a metamodel.
Incremental re-execution after a change of the transformation itself was found only once, assuming an engine that keeps running.
The authors argue for this extension rather than providing a demonstration of it \cite{hearnden2006incremental}.
Vitruvius has the capability to incorporate existing artifacts into an empty V-SUM, but with reactions that do not change \cite{leonhardt2015integration}.
Alternative approaches to consistency treat their own specification as a given \cite{barkowsky2023incremental,konig2025transforming}.
A proposal within Vitruvius aims to establish consistency preservation rules between deltas rather than between metamodels.
Its variants start from a consistent core model and are derived under rules that do not change \cite{neumann2025towards}.
No publication was found that altered the consistency specification of a previously recorded V-SUM, which is the gap this thesis addresses.